\begin{theorem}[核版 RH 临界点先于最近复分歧，且边界为纯谱相变]\label{thm:xi-rh-critical-before-nearest-complex-branch}
\leanverified{paper\_xi\_rh\_critical\_before\_nearest\_complex\_branch}
沿用推论 \ref{cor:real-input-40-collision-branch-radius} 与
命题 \ref{prop:real-input-40-collision-rhk-window} 的记号。
令
\[
P_2(\theta):=\log \rho(e^\theta).
\]
则：
\begin{enumerate}
  \item \(P_2(\theta)\) 的代数分歧点恰为
  \[
  \theta=\log u_j+2\pi i k,\qquad k\in\ZZ,
  \]
  其中 \(u_j\) 遍历三次判别式方程
  \[
  \Delta(u)=4u^3+25u^2+88u+8=0
  \]
  的三个根；
  \item 在 \(\theta=0\) 处的 Taylor 收敛半径为
  \[
  R_\theta=\min_j\min_{k\in\ZZ}\bigl|\log u_j+2\pi i k\bigr|
  \approx 2.76230289346087;
  \]
  \item 核版 RH/Ramanujan 条件的唯一临界参数满足
  \[
  u_R^5+4u_R^4+3u_R^3-96u_R^2-42u_R-14=0,
  \qquad
  \theta_R:=\log u_R\approx 1.27888,
  \]
  且
  \[
  \textsf{RH}(u)\Longleftrightarrow 0<u\le u_R.
  \]
\end{enumerate}
特别地，
\[
\boxed{\;
\theta_R<R_\theta.\;}
\]
并且在临界点 \(u_R\) 处有
\[
\Lambda(u_R)=\lambda_1(u_R)^{1/2},
\]
其中 \(\Lambda(u)\) 为命题 \ref{prop:real-input-40-collision-rhk-window} 的次谱半径。因而存在非空实区间
\[
(\theta_R,R_\theta)\subset \RR_{>0}
\]
使得在该区间上 \(P_2(\theta)\) 仍解析，而 kernel RH 条件已经失效。
换言之，此处 RH 型失效并非由解析边界触发，而是由解析域内部的纯谱交叉触发。
\end{theorem}
\begin{proof}
前三条分别是推论 \ref{cor:real-input-40-collision-branch-radius} 与
命题 \ref{prop:real-input-40-collision-rhk-window} 的直接重述。
故只需比较显式数值：
\[
\theta_R\approx 1.27888
\qquad\text{而}\qquad
R_\theta\approx 2.76230289346087.
\]
于是立即得到
\[
\theta_R<R_\theta.
\]
若 \(\theta\in(\theta_R,R_\theta)\)，则一方面 \(|\theta|<R_\theta\)，故
\(P_2(\theta)\) 仍位于 \(\theta=0\) 的解析收敛盘内；
另一方面 \(u=e^\theta>u_R\)，由
\(
\textsf{RH}(u)\Longleftrightarrow 0<u\le u_R
\)
知 kernel RH 条件已失效。
又由命题 \ref{prop:real-input-40-collision-rhk-window} 的 RH-sharp 等号，边界点满足
\(\Lambda(u_R)=\lambda_1(u_R)^{1/2}\)。故所述“解析仍成立而 RH 已失效”的区间确实非空，并且临界边界由次谱半径与平方根 Perron 尺度的等号刻画，而不是由分支奇点或极点触发。
证毕。
\end{proof}
